\section{Replication Results}

\subsection{Standalone sleeve economics}

\begin{table}[H]
\centering
\caption{Gross strategy performance}
\label{tab:performance}
\begin{tabular}{lrrr}
\toprule
Strategy & Annual return & Sharpe & Maximum drawdown \\
\midrule
Momentum & 13.25\% & 0.46 & $-39.35\%$ \\
Reversal & $-2.89\%$ & $-0.10$ & $-50.48\%$ \\
Static 50/50 & 5.18\% & 0.33 & $-17.61\%$ \\
Primary router & 13.69\% & 0.51 & $-36.24\%$ \\
Equal-proxy router & 7.77\% & 0.30 & $-47.37\%$ \\
\bottomrule
\end{tabular}
\end{table}

Momentum is positive in the evaluation sample, although its risk-adjusted performance is
modest and its drawdown is large. Reversal has negative unconditional economics. This is
already a challenge for the source mechanism: a router can create stable value only if
the replacement sleeve is economically useful in the states in which it is selected.

The primary router has the highest endpoint and a slightly smaller drawdown than
momentum. It does not dominate the static blend on drawdown, and the equal-proxy router
performs materially worse than momentum. \Cref{fig:performance} illustrates why endpoint
ranking alone is an incomplete replication criterion.

\begin{figure}[H]
\centering
\includegraphics[width=0.96\textwidth]{fig01_public_performance.png}
\caption{Gross performance and path risk. The two routers differ only in the definition
of aggregate market volatility.}
\label{fig:performance}
\end{figure}

\subsection{The conditional relation is neither monotonic nor proxy-stable}

Under the primary proxy, reversal exceeds momentum in Q2 and Q4, while momentum strongly
dominates in Q3. In Q4 both sleeves lose money: momentum averages $-0.71\%$ per month and
reversal averages $-0.54\%$. The 17 basis-point difference satisfies the source
direction only weakly and does not show that reversal becomes a positive high-volatility
sleeve.

Under equal-product volatility, the Q4 ranking reverses. Momentum averages 2.40\% per
month and reversal 0.90\%. Q3, not Q4, is the state in which reversal exceeds momentum.
The conditional pattern is therefore not a monotonic function of volatility and is not
stable to a reasonable market definition.

\begin{figure}[H]
\centering
\includegraphics[width=0.96\textwidth]{fig02_public_conditional_returns.png}
\caption{Conditional gross returns by the lagged volatility state used for the holding
month.}
\label{fig:conditional}
\end{figure}

\subsection{The primary improvement is event-concentrated}

The primary router improves on momentum by only 3.69 basis points per month on average,
with a HAC $t$-statistic of 0.06. Among 14 Q4 months, reversal wins four. Its mean
advantage over momentum is 0.17\% per month, but the median is $-2.68\%$. One event has a
34.28 percentage-point reversal advantage; excluding it, the mean Q4 advantage is
$-2.45\%$ per month.

\begin{figure}[H]
\centering
\includegraphics[width=0.78\textwidth]{fig05_public_event_influence.png}
\caption{Aggregate Q4 influence diagnostics. The event date and return series are not
included in the sanitized evidence package.}
\label{fig:influence}
\end{figure}

This pattern is better described as an event-level momentum-crash hedge echo than a
stable Q4 routing relation. A single successful rescue is economically interesting, but
it does not make high volatility a sufficient switching condition.
